\documentclass[11pt]{article}
\usepackage[T1]{fontenc}
\usepackage[utf8]{inputenc}
\usepackage{lmodern}
\usepackage[margin=1in]{geometry}
\usepackage{amsmath,amssymb}
\usepackage[numbers]{natbib}
\usepackage[hidelinks]{hyperref}
\title{Self Without a Transcript:\\Semantic Recurrence Across Text and Latent States}
\author{}
\date{November 1, 2026}
\begin{document}
\maketitle
\begin{abstract}
We study what changes when a frozen language model verbalizes a hidden state and re-encodes its own words. Text relay (T), latent recurrence (L) and a language roundtrip after the same latent update (B) each retain only the current state. In four Mistral-7B-Instruct-v0.3 development trials, we varied the stopping rule while holding two initial descriptions and the layer boundaries fixed. All 18 generations in the EOS condition ended naturally. B produced one stored state across the four updates under each stopping rule, while L produced four. We examine whether a separate readout can distinguish sources even when greedy verbalization coincides. The latent transport is uncalibrated and held-out task transfer remains open.
\end{abstract}

\section{Problem formulation}
A description such as ``I prioritize keeping decisions revisable'' may be reinterpreted as a preference for preserving correction opportunities. The transition may add a useful distinction, remain a paraphrase, or introduce an unsupported claim. Our object is the change in meaning, not the event of uttering the earlier description. We exclude a direct transcript and instructions that infer character from the act of speaking; this does not erase past updates' causal influence on the current state.

We ask whether recurrence yields defensible distinctions, whether a language roundtrip changes their formation or use, and whether changes improve held-out problem solving or appropriate switching under reversed conditions. Convergence is not intrinsically failure, and novelty is not intrinsically success. ``Self'' denotes an attribution frame rather than a claim about subjective experience.

\section{Related work}
Coconut trains continuous-state reasoning through hidden-state feedback \citep{hao2024coconut}. Recurrent-depth models study repeated internal computation \citep{geiping2025huginn}; STARS motivates distinguishing useful stable dynamics from collapse \citep{yang2026stars}. These do not establish the validity of arbitrary untrained residual feedback.

Persona vectors and the Assistant Axis connect persona-related behavior with activation structure and interventions \citep{chen2025persona,lu2026axis}. Natural Language Autoencoders provide a learned activation/text roundtrip \citep{frasertaliente2026nla}, while latent chain-of-thought analysis distinguishes output and representational commitment \citep{li2026dynamics}. Our estimand concerns a specific codec rather than ``language versus pure thought.''

Semantic diversity differs from surface variation \citep{tao2026divergent}. Contrasting self-correction and condition-verification findings motivate operation-specific hypotheses \citep{huang2024selfcorrect,wu2024verification}. Dadfar's report is considered with McCardle's open replication critique concerning length, mode and cycle confounds \citep{dadfar2026examine,mccardle2026replication}. SlimMoE informs the proposed Phi MoE comparison \citep{li2025slimmoe}. We claim neither priority nor a systematic review.

\section{Three conditions}
Let $c_n$ be a fixed-schedule operation instruction, $s_n$ a current description, and $Z_n\in\mathbb R^{K\times d}$ a state. With encoder $E$, verbalizer $D$, and state updater $G$:
\begin{align}
\text{T:}\quad s_{n+1}&=T_\theta(c_n,s_n),\\
\text{L:}\quad Z_{n+1}&=G_\theta(c_n,Z_n),\\
\text{B:}\quad Z_{n+1}&=E_\theta\bigl(D_\theta(G_\theta(c_n,Z_n))\bigr).
\end{align}
L and B begin from copies of $E(s_0)$ and share weights, boundaries, transport, generation policy and schedule. Their contrast is primary; T/L is a system-level comparison.

Each outer step uses a fresh context without persistent KV cache. Only current text or slots carry forward. Observation uses clones and restores RNG state; its text is not fed back except as prescribed by B. The initial description is not repeatedly reintroduced. Slots occupy explicit decoder boundaries; post-writer to pre-cut transport is explicit. An identity map is an unqualified engineering diagnostic.

The two schedules are an ordinary meaning update and a fixed sequence of decomposition, boundary examination and reintegration. Unchanged content is allowed when no justified update exists. Adaptive cycle-triggered feedback would be a separate condition.

\section{Qualification and measurement}
Native parity, content discrimination, contrastive-state swap effects, one-step validity, readout fidelity and positional controls precede any latent pilot. Transport or codec training may be needed; training material must remain disjoint from confirmation. One-step success does not establish long-trajectory validity.

\paragraph{Meaning.} Blind paired judgments assess support, update type, new distinctions and event-recovery violations. Lexical and tensor distances are not semantic scores.
\paragraph{Activations.} Selected residual, attention, MLP and supported MoE gate outputs are retained. Gate logits are not actual dispatch frequencies. Independent probes and causal patching remain planned.
\paragraph{Effects.} Problems arrive only after recurrence and operate on isolated copies. Near, structural and less-related transfer are separate tiers. Initial-state, one-shot instruction and non-self controls are planned. Reversals distinguish commitment from rigidity. Frozen-weight state-conditioned gains are not permanent learning.

At a common state, compare $G^r(Z)$ with $G^r(E(D(Z)))$. Continuous bottlenecks, norm/distance-matched perturbations, no-update roundtrips, and meaning-preserving versus contrastive swaps are needed to distinguish codec damage from more specific effects. Only the identity/language branch pair is currently implemented.

\section{Proposed pilot and analysis}
Six content groups, two poles, two paraphrases and self/other frames yield 48 development seeds. Three arms and two schedules yield 288 proposed trajectories per model. Candidate families are Mistral-7B-Instruct-v0.3, Llama-3.2-3B-Instruct, Gemma-2-2B-IT and Phi-mini-MoE, with Phi-tiny-MoE a separately reported fallback. Family differences are not controlled architectural interventions.

Twelve steps with fixed assessment checkpoints are proposed. The present CLI scores tasks at baseline/final only. Its 20 synthetic tasks are exposed development fixtures, not a broad benchmark. Confirmation data, thresholds, sample size and preregistration are pending.

Paired L/B utility differences at a declared depth will be reported alongside support and switching. Inference must account for dependence among meaning groups and task families rather than treating each step as independent. Same-step designs report actual compute; matched-compute analyses are secondary. Failures, censored generations and ties retain denominators. Best-depth selection on confirmation is prohibited.

\section{Development experiments}
\subsection{Native EOS and the language roundtrip}
Two opposed descriptions emphasized revisable or prompt decisions. Each underwent two plain updates in T, L and B, once with native EOS and once with an eight-token ceiling. The four trials used greedy BF16 Mistral-7B-Instruct-v0.3, eight slots, a 1024-token context, a pre-layer-16 cut and a post-layer-31 writer (zero-indexed). All 18 generations per policy include main-path, observer and branch readouts. All 18 reached EOS with the full context available, and all 18 reached the explicit ceiling under the cap. With EOS, each of the four B updates verbalized the same 13-token sentence about workspace contents. The four B states matched within each policy, although the two policies gave different states. The four L states differed within each policy and matched step for step across policies.

\section{Limitations}
These are development measurements on frozen checkpoints. The identity transport and verbalizer have not passed a separately reviewed semantic gate. The source of an observed change can include the reconstruction error, layer coordinates, positions, prompt wording or additional computation. Interpretation of novel distinctions and downstream utility requires independent judgments and confirmation tasks.

\bibliographystyle{plainnat}
\bibliography{references}
\end{document}
